Nous avons utilisé des noms pour désigner des valeurs~: les constantes, les paramètres des
procédures et les compteurs des boucles ``for''. Une fois fixée, la valeur désignée par un tel nom
ne change plus. Les {\em variables} permettent au contraire de mémoriser des résultats
intermédiaires dans les calculs et de les modifier au cours de l'exécution du programme.

\section{Variables et affectation}
Une variable\index{Variable} est utilisée pour conserver une valeur d'un certain type dans la mémoire
de l'ordinateur. En \ocaml, une variable s'appelle une {\em référence}. Une déclaration de
variable\index{Déclaration de variable} indique le nom de la variable (qui est un identificateur)
et sa valeur initiale, placée après le mot {\idt ref}~:

\vspace{\neglskip}\mbox{\blc  \blk\idt  \mcl let\idt \ b\ =\ ref\ true}\par
\vspace{\neglskip}\mbox{\blc  \blk\idt  \mcl let\idt \ i\ =\ ref\ 0}\par
\vspace{\neglskip}\mbox{\blc  \blk\idt  \mcl let\idt \ r\ =\ ref\ 0.0}\par
\vspace{\neglskip}\mbox{\blc  \blk\idt  \mcl let\idt \ s\ =\ ref\ {\cha \char34\char34}}\par

\noindent C'est le compilateur qui décide quelle partie de la mémoire et quel codage sera utilisé pour
représenter la valeur de la variable. Le type de la variable est celui de sa valeur initiale~:
{\idt b} est une variable booléenne, {\idt i} une variable entière, {\idt r} une variable
rationnelle et {\idt s} une variable contenant une chaîne de caractères. En \ocaml, une variable
a donc toujours une valeur, fixée dès sa création. On peut utiliser cette valeur dans une
expression (comme {\idt 2 * !i + 1}) en la notant par le nom de la variable précédé d'un point
d'exclamation ({\idt !i} dans l'exemple, qui se lit ``contenu de {\idt i}''). Il ne faut pas
l'oublier~: si l'on écrit {\idt i + 1} au lieu de {\idt !i + 1}, le compilateur signale
`{\tt The value i has type int ref but an expression was expected of type int}' (la variable
{\idt i} est une référence vers un entier, pas un entier). L'instruction
d'{\em affectation\/} \index{Affectation}permet de changer la
valeur d'une variable. Par exemple l'instruction~:

\vspace{\neglskip}\mbox{\blc  \blk\blk\idt  b\ :=\ false}\par

\noindent affecte la valeur ``faux'' à la variable booléenne {\idt b}. De manière générale, une
instruction d'affectation a la forme suivante~:

\index{Affectation}
\begin{center}
\begin{tabular}{@{}|p{12.5cm}|@{}} \hline
{\em Instruction d'affectation~:}\\ \hline \hline
\begin{minipage}[t]{12.5cm}
{\medskip

\vspace{\neglskip}\mbox{\blc  \blk\blk\idt  nomDeVariable\ :=\ \rm expression%
\footnote{\ Le type de l'expression est celui de la valeur initiale de la variable.}}\par

\medskip}
\end{minipage}\\ \hline
\end{tabular}
\end{center}

\noindent Le symbole\ {\idt :=}\ se lit ``devient égal à''. L'exécution de cette instruction consiste à
évaluer l'``expression'' pour calculer sa valeur qui devient celle de la variable nommée
{\idt  nomDeVariable}. Le type de l'expression doit être exactement celui de la variable~:

\index{Compatibilité de type}
\begin{center}
\begin{minipage}[t]{11.5cm}
\begin{tabular}{@{}|l|l|l|@{}} \hline
\multicolumn{3}{@{}|l|@{}}{\em Types pour l'affectation~:} \\ \hline \hline
Déclaration & Type de la variable & Type de l'expression\\ \hline
{\idt ref true} & {\idt bool ref} & {\idt bool} \\
{\idt ref 0} & {\idt int ref} & {\idt int} \\
{\idt ref 0.0} & {\idt float ref} & {\idt float} \\
{\idt ref {\cha \char34\char34}} & {\idt string ref} & {\idt string} \\
\hline
\end{tabular}
\end{minipage}
\end{center}

\noindent  Par exemple, si {\idt i} est une variable entière, l'affectation {\idt i := 1.5} est
refusée par le compilateur qui signale `{\tt The constant 1.5 has type float but an expression was
expected of type int}'. Comme pour les paramètres, le compilateur ne fait jamais de
{\em conversion}\index{Conversion} de lui-même~: il faut écrire par exemple
{\idt r := float\_of\_int !i} ou {\idt i := truncate !r}.

\section{Sauvegarde de l'état du robot}

Une première utilisation des noms est la sauvegarde de l'état du \robot. Imaginons par
exemple que nous voulions concevoir une procédure pour faire un dessin qui nécessite divers
changements de la taille de la grille. Pour que l'appel de
cette procédure dans un programme ne modifie pas la taille de la grille utilisée dans le
programme, il faut sauvegarder cette
taille au début de la procédure puis la restituer à la fin. Comme la taille sauvegardée ne change
pas pendant le dessin, il n'est pas nécessaire d'utiliser une variable~: il suffit de lui donner un
nom par une {\em déclaration locale}\index{Déclaration locale}
{\mcl let}\ {\idt nom = expression} {\mcl in}, comme dans l'exemple suivant~:

{\raggedright
 %\vspace{\baselineskip}
 \vbox to 0pt{\includevspacePostScript[13mm]{Canards}{0.95}{0.95}}%
 \input ocaml/Canards.tex
 }

\noindent La déclaration {\mcl let}\ {\idt lX = valLgX ()} {\mcl and} {\idt lY = valLgY ()} {\mcl in}
calcule les valeurs de {\idt valLgX ()} et de {\idt valLgY ()} au début de la procédure et
leur donne les noms {\idt lX} et {\idt lY}, qui ne peuvent être utilisés que dans la suite du
corps de la procédure {\idt canard} (après le mot {\mcl in}).

\section{Compter}
Supposons que nous voulions
compter le nombre d'itérations dans une boucle ``while''. On utilise une variable entière~:

\mbox{\blc  \blk\idt  \mcl let\idt\ i\ =\ ref\ 0}\par

\noindent Avant d'exécuter la boucle ``while'', la variable doit être initialisée à 0
(elle peut avoir été modifiée depuis sa déclaration)~:

\mbox{\blc  \blk\blk\idt  i\ :=\ 0}\par

\noindent Après chaque tour de boucle, la valeur de la variable doit être augmentée de 1~:

\mbox{\blc  \blk\blk\idt  i\ :=\ !i\ +\ 1}\par

\noindent Supposons que la valeur de la variable {\idt i} soit égale à 5 {\em avant} l'affectation
\mbox{\idt  i\ :=\ !i\ +\ 1}. La valeur de l'expression {\idt  !i\ +\ 1}
est $5 + 1 = 6$ de sorte que la valeur de la variable {\idt i} {\em après} l'affectation sera égale
à 6.

Par exemple, ceci peut nous servir pour déterminer la taille de la fenêtre dans laquelle le \robot\
se déplace sur l'écran de l'ordinateur~:
\input ocaml/CalculerTailleFenetreRobot.tex

\noindent L'opérateur {\idt \char94} met bout à bout deux chaînes de caractères et la
fonction {\idt string\_of\_int} donne la chaîne de caractères qui représente un entier (par
exemple {\idt string\_of\_int 25} vaut {\cha \char34{}25\char34}).

\section{Entrée d'une valeur au clavier}
On peut obtenir une valeur en la tapant au clavier de
l'ordinateur. Pour
ce faire la fonction {\idt lireEntier ()} attend que l'on tape un nombre entier et rend
sa valeur. La question posée par le programme (écrite avec {\idt write}, voir ci-dessous)
s'affiche dans la barre du bas de la fenêtre du \robot, et la réponse que l'on tape au clavier
s'inscrit à sa suite (on peut aussi la taper dans le terminal d'où le programme a été lancé).
On écrit par exemple~:

\mbox{\blc  \blk\blk\idt  write\ [{\cha \char34{}n\ =\ \char34}];}\par
\mbox{\blc  \blk\blk\idt  \mcl let\idt \ n\ =\ lireEntier\ ()\ \mcl in}\par

\noindent pour donner au nom {\idt n} la valeur tapée, ou {\idt i := lireEntier ()} pour
l'affecter à une variable {\idt i}. On peut
corriger les fautes de frappe avec la touche  %touche d'effacement
\framebox[11.mm][l]{$\leftarrow$~~~}\hspace{0.3ex}
d'effacement. Il faut terminer l'entrée de la valeur par la frappe de la touche %retour a la ligne
\framebox[11.mm][l]{$\hookleftarrow$~~~}\hspace{0.3ex}
de retour à la ligne. Si ce que l'on a tapé n'est pas un entier, le programme le
redemande. Sur un clavier français, on peut taper les chiffres avec ou sans la touche
des majuscules.

De même {\idt lireReel ()} rend un nombre rationnel (de type {\idt float}), {\idt lireChaine ()}
rend la chaîne de caractères tapée (sans apostrophes ni guillemets) et {\idt lireCar ()} rend
le caractère de la première touche frappée dans la fenêtre du \robot\ (sans attendre la touche de
retour à la ligne). En \pascal, une seule procédure {\idt readln} permettait de lire une valeur
de n'importe quel type. En \ocaml, chaque fonction de lecture a un type de résultat bien précis, ce
qui permet au compilateur de vérifier que la valeur lue est utilisée correctement. Certaines
fonctions sont cependant {\em polymorphes}\index{Polymorphe}, comme {\idt ref} qui accepte
une valeur initiale de n'importe quel type (alors que la plupart des fonctions, comme
{\idt lireEntier} ou {\idt sqrt}, sont {\em monomorphes}\index{Monomorphe} en ce sens que leurs
paramètres et leur résultat ont toujours le même type).

L'instruction {\idt readln ()} permet d'attendre que la touche %retour a la ligne
\framebox[11.mm][l]{$\hookleftarrow$~~~}\hspace{0.3ex}
de retour à la ligne soit enfoncée. Elle est donc utile pour laisser le temps de lire ce que
le programme a écrit avant la fin de son exécution.

\section{Affichage de la valeur d'une expression}
Pour afficher des textes dans le terminal (et dans la barre du bas de la fenêtre du
\robot), on
peut utiliser l'instruction {\idt write [}$t_1${\idt ;} $t_2${\idt ;} \ldots{\idt ]} qui écrit les chaînes de
caractères $t_1$, $t_2$, \ldots\ les unes
à la suite des autres, sur la même ligne. L'instruction
{\idt writeln [}$t_1${\idt ;} $t_2${\idt ;} \ldots{\idt ]} fait de même puis passe à la ligne suivante. Pour
écrire la valeur d'une expression entière $E$, on la transforme en chaîne de caractères
avec {\idt string\_of\_int} $E$ (ou {\idt string\_of\_float} $E$ pour une expression
rationnelle). Par exemple
le programme ci-dessous lit
un entier $n$
puis écrit sa factorielle $n! = (1 \ast 2 \ast \ldots \ast (n - 1) \ast n)$. Pour redemander
la valeur de $n$ tant qu'elle est incorrecte, il utilise une fonction récursive {\idt lireN}
(voir le chapitre \ref{ch:rec}) qui rend la valeur lue~:
\input ocaml/Factorielle.tex

\noindent La valeur maximale de $n$ est 20 car $21!$ dépasse le plus grand entier
{\idt max\_int} $= 4\,611\,686\,018\,427\,387\,903$. Au-delà, le calcul serait faux sans
que l'ordinateur le signale (par exemple {\idt max\_int + 1} vaut {\idt min\_int}, le plus petit
entier négatif$\:$!).

\vfill

\noindent \rule{137mm}{0.1mm} %---------------------------------------------

\noindent\begin{exercices}{\ \\

\tiret Écrire un programme qui dessine un cadre le plus près possible du bord de la fenêtre de déplacement
du \robot.

\tiret Écrire un programme qui esquisse les dessins à  compléter de la figure \ref{fig:completer-les-dessins} 
(page \pageref{fig:completer-les-dessins}).

\tiret La numération babylonienne \cite{guitel}, qui date d'environ 2000 ans av. J.-C., est en
base 60. C'est donc une numération sexagésimale tout comme celle que nous utilisons encore
pour noter l'heure. Les unités, de 1 à 59, sont notées de manière additive en accumulant des ``clous'' (pour
1) et des ``chevrons'' (pour 10) regroupés et superposés comme indiqué sur la figure
ci-dessous.\\[3.9cm]
\noindent\includePostScript{Babyloniens}{0.95}{0.95}%[3.9cm]

\noindent Au delà de 60, la notation est positionnelle, la valeur d'un chiffre occupant
la $p^{\grave{e}me}$ position (comptée de droite à gauche à partir de zéro) étant
multipliée par $60^{p}$. L'absence d'unités sexagésimales en position médiale (que
nous no\-te\-rions par 0) fut d'abord marquée par un espace puis, vers le $3^{\grave{e}me}$
siècle av. J.-C., par ``deux clous inclinés''. Les Babyloniens ont donc mis 15 siècles pour
inventer une notation pour le
zéro sans toutefois en comprendre la valeur calculatoire$\:$! Par exemple 
le nombre $3657 = 1 \times 60^{2} + 0 \times 60^{1} + 57 \times 60^{0}$
est représenté par la juxtaposition des chiffres $1$, $0$ et 
$57 = 50 + 7$~:

\pagebreak
%\vspace{-6mm}
\mbox{}\\[1.9cm]
\noindent\includePostScript{Babyloniens3657}{0.95}{0.95}%[2.4cm]

Écrire un programme \ocaml\ de conversion
d'un nombre po\-si\-tif quel\-con\-que dans la nu\-mé\-ra\-tion babylonienne. 

\tiret La numération égyptienne ($3^{\grave{e}me}$ 
millénaire av. J.-C.) est décimale et additive \cite{guitel}. Elle utilise des hiéroglyphes pour 
noter les
puissances de 10 (le trait vertical pour 1, l'anse pour 10, la spirale pour 100, le lotus pour 1000,
l'index pour 10~000,
le têtard pour 100 000 et le génie pour 1 000 000)~:\\[3.0cm]

\noindent\includePostScript{ChiffresEgyptiens}{0.95}{0.95}%[3.0cm]
 
\noindent La valeur d'une séquence de signes est la somme des valeurs
de ces signes~:\\[2.0cm]

\noindent\includePostScript{1234567Egyptien}{0.95}{0.95}%[3.0cm]
 
\noindent On écrit de gauche à droite ou de droite à gauche, la
lecture se faisant dans le sens où sont orientés les hiéroglyphes. 
Les signes identiques se regroupent sur une, deux ou trois lignes de
la manière suivante~:\\[0.8cm]
 
%\vspace{-\baselineskip}
\noindent\includePostScript{GroupesEgyptiens}{0.95}{0.95}%[1.8cm]
 
\noindent  Entre -2000 et -1600 ans av. J.-C., une représentation hybride (additive et partiellement
positionnelle) est apparue. Un nombre inférieur à 1000 situé avant, au dessus ou au dessous
du signe du génie du million, du têtard, de l'index ou du lotus représente le produit de ce nombre et de la
valeur de ce hiéroglyphe. Par exemple, le nombre 37 280 549 peut se décomposer en
$((372 \ast 100 000) + 80 549)$ et s'écrit~:\\[2.3cm]
 
\noindent\includePostScript{37280549Egyptien}{0.95}{0.95}
 
\noindent Écrire un programme \ocaml\ de conversion d'un nombre dans la nu\-mé\-ra\-tion
é\-gyp\-tien\-ne.

\tiret La numération acrophonique grecque \cite{guitel} fut utilisée à partir 
du V$\grave{e}me$
siècle av. J.-C. à Athènes. C'est une numération additive décimale utilisant la base
auxiliaire 5 pour diminuer les répétitions. Les chiffres supérieurs
à un sont notés par des 
lettres (ou combinaisons de lettres) 
correspondant chacune à l'initiale (ou à l'abréviation) du nom du
nombre correspondant~:\\[1.8cm]

\noindent\includePostScript{ChiffresGrecques}{0.95}{0.95}

Écrire un programme \ocaml\ de conversion d'un nombre in\-fé\-rieur à 100~000 dans la nu\-mé\-ra\-tion
grecque. 
Pour représenter les grands nombres, utiliser la numération acrophonique 
pour les monnaies~:

\noindent Oboles~:\\[0mm]

%\vspace{-0.5cm}
\noindent\includePostScript{ObolesGrecques}{0.95}{0.95}

\noindent Drachmes (1 drachme = 6 oboles)~:\\[0.8cm]

\noindent\includePostScript{DrachmesGrecques}{0.95}{0.95}

\noindent Talents (1 talent = 6~000 drachmes)~:\\[0.8cm]

\noindent\includePostScript{TalentsGrecques}{0.95}{0.95}

\clearpage
\noindent Par exemple le nombre 239~686~966 se représente par~:\\[1.0cm]

\noindent\includePostScript{239686966Grecque}{0.95}{0.95}

\tiret La numération romaine \cite{guitel} est additive et soustractive (tout signe placé à gauche
d'un signe de valeur supérieure s'en soustrait). Ses chiffres
sont bien connus~:\\[1.2cm]

\noindent\includePostScript{ChiffresRomains}{0.95}{0.95}

Pour les grands nombres, la barre horizontale multiplie par 1~000 la valeur du nombre qu'elle
surligne et le rectangle sans base multiplie par 100~000 la valeur du nombre qu'il enferme,
comme dans l'exemple 135~792~468 ci-dessous~:\\[1.4cm]

\noindent\includePostScript{135792468Romain}{0.95}{0.95}

Écrire un programme \ocaml\ de conversion d'un nombre dans la nu\-mé\-ra\-tion
ro\-mai\-ne.

\tiret Écrire un programme de dessin de la courbe de Sierpinski à un ordre
quelconque de la figure \ref{fig:CourbSierp} (page \pageref{fig:CourbSierp}).

\tiret Un carré magique d'ordre $n$ est un tableau carré de $n^{2}$ nombres
tel que la somme des nombres de chaque ligne, de chaque colonne et de chaque
dia\-go\-na\-le est la même. Des exemples sont donnés à la figure \ref{fig:Cmagic}
(page \pageref{fig:Cmagic}). Notons $C_{i, j}$ le nombre figurant à la
$i^{\grave{e}me}$ ligne et la $j^{\grave{e}me}$ colonne du carré comptées à
partir de 1 et numérotées de bas en haut pour les lignes et de gauche à
droite pour les colonnes. Écrire un programme \ocaml\ qui cons\-truit des
carrés magiques d'ordre $n$ impair en commençant par placer 1 dans
$C_{i, j} = C_{\frac{n + 1}{2}, n}$ puis affecte les nombres 2, 3, \ldots, $n^{2}$
aux éléments $C_{i, j}$ obtenus en augmentant $i$ et $j$ de 1 modulo $n$ à chaque fois
pendant $n - 1$ pas, et en diminuant $j$ de 1 tout en laissant $i$ inchangé à chaque 
$n^{\grave{e}me}$ pas.

\noindent $\Box$}\end{exercices}
\begin{figure}[htbp]
 \includevspacePostScript[9.2cm]{CourbeSierpinski0a5}{0.95}{0.95}
 \caption{Courbes de Sierpinski d'ordres 0 à 5}\label{fig:CourbSierp}
\end{figure}
\begin{figure}[htbp]
 \includevspacePostScript[6.0cm]{CarreMagique}{0.95}{0.95}
 \caption{Carrés magiques d'ordres 3, 5 et 7}\label{fig:Cmagic}
\end{figure}

\pagebreak

\begin{corrige}{
Nous donnons un programme (sans le corps de la procédure technique {\idt ecrireCentre}).
Les autres se trouvent dans le dossier\newline {\idt livre/\allowbreak 22\_variables\_E\_S}~:
\input ocaml/CarreMagique.tex
\noindent $\Box$}\end{corrige}

%%%%%%%%%%%%%%%%%%%%%%%%%%%%%%%%%%%%%%%%%%%%%%%%%%%%%%%%%%%%%%%%%%%%%%%%%%%
\noindent \framebox[\textwidth]{
\begin{minipage}{\myparboxwidth}
\setlength{\parindent}{2em}
\small \sf
\vspace{\abovedisplayskip}
\section*{Addenda}
Une \underline{variable}\index{Variable} sert à donner un {\em nom} à une mémoire qui conserve
une {\em valeur} représentée
selon un code déterminé par le {\em type} de la variable. Pour un codage binaire, on
peut imaginer que la mémoire est une suite de lampes éteintes (pour représenter
le bit 0) ou allumées (pour le bit 1).
Dans beaucoup de langages, comme \pascal, la valeur
initiale d'une va\-ria\-ble est {\em indéterminée} puisqu'elle dépend de l'état
de la mémoire au moment du début de l'exécution du programme, ce qui est une source
d'erreurs. En \ocaml, une variable reçoit une valeur initiale dès sa création .
L'instruction d'\underline{affectation}\index{Affectation} permet de
conserver une nouvelle valeur en mo\-di\-fiant physiquement
l'état de cette mémoire. La nouvelle valeur est donnée par une expression (qui peut
utiliser l'ancienne valeur). La notion informatique de variable est donc très différente
de la notion mathématique homonyme. En mathématiques une variable désigne une valeur
inconnue fixée une fois pour toutes, comme les noms introduits par {\mcl let} en
\ocaml. Pour un mathématicien, une variable informatique
serait la suite des valeurs qu'elle prend à sa création puis après chaque
affectation.

\vspace{\baselineskip}

On peut déclarer des variables {\em locales} à une procédure (ou fonction),
en écrivant {\mcl let}\ {\idt x = ref} $\ldots$ {\mcl in} dans le corps de la procédure, comme
les variables {\idt i} et {\idt j} du programme {\idt CarreMagique}.
Dans ce cas la variable n'est utilisable que dans la suite du corps de la procédure.
À l'exécution, elle est créée à chaque exécution du corps
de la procédure. Pour une procédure récursive,
il s'agit à chaque appel d'une nouvelle variable sans aucun
lien avec les instances correspondant aux appels récursifs antérieurs.

Dans le corps de la procédure on peut également utiliser des variables
{\em glo\-ba\-les} dont la déclaration figure dans le programme
avant celle de la procédure, comme {\idt largeur} et {\idt hauteur} dans le programme
{\idt CalculerTailleFenetreRobot}. Ces variables globales existent
pendant toute la durée de l'exécution du programme. Par conséquent en cas
d'appels récursifs d'une procédure, il s'agit toujours de la
même variable globale.

Enfin des variables locales et globales peuvent porter le même nom. Dans
ce cas c'est la variable locale qui est visible dans la procédure, la variable
globale continuant à exister mais étant inaccessible.

\vspace{\abovedisplayskip}
\end{minipage}
}
%%%%%%%%%%%%%%%%%%%%%%%%%%%%%%%%%%%%%%%%%%%%%%%%%%%%%%%%%%%%%%%%%%%%%%%%%%%% 

